\begin{description}[leftmargin=1.5em, style=nextline, font=\bfseries]
    \item[$\top$] 
        Rappresenta la verità logica, la sua prova è costituita da una costruzione canonica banale che è sempre disponibile 
        e non richiede informazioni preresse.
    \item[$\bot$] 
        Rappresenta l'assurdo logico e coerentemente non possiede alcuna prova, non è possibile fornire alcuna 
        costruzione valida per questo enunciato.
    \item[$\phi \land \psi$] 
        Una prova è costituita da una coppia ordinata $\langle e_1, e_2 \rangle$ che contiene simultaneamente un'
        evidenza $e_1$ per $\phi$ e un'evidenza $e_2$ per $\psi$.
    \item[$\phi \lor \psi$] 
        Una prova è una coppia $\langle i, e \rangle$ dove l'indice $i \in \{0, 1\}$ agisce da selettore: se $i=0$ 
        l'evidenza $e$ prova il disgiunto $\phi$, mentre se $i=1$ essa prova il disgiunto $\psi$.
    \item[$\phi \to \psi$] 
        Una prova è una costruzione funzionale $f$ in grado di trasformare ogni possibile prova di $\phi$ in una prova 
        di $\psi$, definendo così un metodo effettivo di derivazione.
    \item[$\neg \phi$] 
        Equivalente a $\phi \to \bot$, la sua prova è una costruzione che trasforma ogni ipotetica evidenza di 
        $\phi$ in una prova della contraddizione ($\bot$), mostrandone l'asseribilità impossibile.
\end{description}